\subsection{Uniqueness for the constructed torus flow}
The residual identity gives the solution. We now show that its particle
event is determined by the force and initial data. The proof is the
classical difference-energy comparison \cite[Section~18]{Leray}.
If $(v,p)$ is another solution in the classical class of
Section~\ref{sec:statement}, set $w=v-U$. Subtraction gives
\[
 w_t+(v\cdot\nabla)w+(w\cdot\nabla)U
  =-\nabla p+\nu\Delta w,\qquad \diver w=0.
\]
Multiplication by $w$ and periodic integration by parts yield
\[
 \frac12\frac{d}{dt}\|w\|_2^2+\nu\|\nabla w\|_2^2
 \leq\|\nabla U\|_{\infty,\mathrm{op}}\|w\|_2^2.
\]
The stated continuity of $v_t$, the first two spatial derivatives and
pressure gradient justifies all these operations. The transport and
pressure terms vanish by incompressibility. Boundedness of $\nabla U$
on each finite interval and $w(0)=0$ give $w=0$ by Gronwall's
inequality. The equation then gives $\nabla p=0$, and mean-zero
normalization gives $p=0$. Smoothness and bounded spatial gradient on
the compact torus give a unique material flow through every finite
time. Thus the event already proved for $U$ belongs to the unique
fluid solution. Appendix~\ref{app:comparison} supplies the corresponding
noncompact comparison theorem and the finite-time flow identities.

\subsection{Solenoidal forcing with the same velocity}
The Helmholtz--Leray projection removes the divergence of the force.
Because it changes a gradient, it also changes the pressure. The
following quantitative version is useful both for the present
$L^2$ bounds and for other time profiles.

\begin{lemma}[Effective periodic projection]\label{lem:projection}
Let $g$ be a smooth mean-zero vector field on
$[0,\infty)\times\T^3$, effectively evaluable together with all
mixed derivatives and with effective bounds for those derivatives
on every finite time interval. There is a unique smooth mean-zero
$\phi$ satisfying $\Delta\phi=\diver g$. Both $\phi$ and
$g_{\rm sol}=g-\nabla\phi$ are effective, and $g_{\rm sol}$ is
divergence free and mean zero. For every $h\geq0$ and spatial
multi-index $\mu$,
\begin{equation}\label{eq:projection-estimate}
 \|\partial_t^h\partial_X^\mu\nabla\phi(t)\|_\infty
 \leq C_\mu\max_{1\leq i\leq3}
 \|\partial_t^h\partial_{X_i}^{|\mu|+5}g(t)\|_\infty.
\end{equation}
Thus, when all mixed derivatives of $g$ are bounded, square-integrable
in time in spatial supremum norm, or rapidly decaying, the same
property holds for $g_{\rm sol}$. Here rapid decay means a finite
supremum after multiplication by $(1+t)^A$, for every $A\geq0$ and
every fixed derivative order. Time periodicity and stationarity
after a given time are also preserved. If $(U,P)$ solves the equation
with force $g$, then it solves it with force $g_{\rm sol}$ and
pressure $P-\phi$.
More precisely, any common pointwise time weight on the finitely many
derivatives on the right of \eqref{eq:projection-estimate} passes to
the derivative on its left. Thus a specified polynomial bound is
preserved whenever those higher spatial derivatives have that bound.
\end{lemma}
\begin{proof}
For the Fourier basis $e^{2\pi i k\cdot X}$ set
\[
 \widehat\phi(k)=-\frac{i k\cdot\widehat g(k)}{2\pi|k|^2}
 \quad(k\ne0),\qquad \widehat\phi(0)=0.
\]
The multiplier for $\nabla\phi$ is $kk^{\mathsf T}/|k|^2$,
whose operator norm is one. For a derivative of order $r$, integrate
the coefficients of $\partial_t^hg$ by parts $r+5$ times in a
coordinate with maximal $|k_i|$. A term of the differentiated series
is bounded by a fixed multiple of $|k|^{-5}$ times the right-hand
derivative norm in \eqref{eq:projection-estimate}. There are
$O(n^2)$ points in the shell $\|k\|_\infty=n$. The series and its
derivatives therefore converge absolutely, with an explicit tail
bounded by $C\sum_{n>K}n^{-3}$. The series for $\phi$ has one
further factor $|k|^{-1}$. Termwise differentiation proves the
Poisson equation; the Fourier coefficients also prove uniqueness
among mean-zero solutions and establish
\eqref{eq:projection-estimate}.

The Fourier identification used here is justified in
Appendix~\ref{app:fourier}. Since $g$ is real,
$\widehat\phi(-k)=\overline{\widehat\phi(k)}$, so $\phi$ is real.

Riemann sums with a bound on one more spatial derivative compute
each coefficient to any requested error. Combined with the tail
estimate, this gives effective evaluation of every derivative.
The estimate is pointwise in time and commutes with time
derivatives, so it proves each stated norm and weighted-decay
assertion. Linearity and uniqueness of the mean-zero Poisson
solution preserve periodicity and eventual stationarity. Finally,
\[
 -\nabla(P-\phi)+(g-\nabla\phi)=-\nabla P+g,
\]
which proves the pressure formula.
Conversely, a classical solution with force $g_{\rm sol}$ and pressure
$p$ solves the equation with force $g$ and pressure $p+\phi$.
This correspondence preserves the normalized classical class on the
torus, since $\phi$ is smooth and mean zero. It therefore preserves
uniqueness for fixed initial data, as well as the reference velocity.
\end{proof}

To apply Lemma~\ref{lem:projection} to the force in \eqref{eq:force},
we still need to verify its effective-evaluation hypotheses. The next
subsection supplies them and completes the optional solenoidal
construction. This periodic projection does not assert preservation
of compact support on a Euclidean domain, nor does it preserve
pressure zero.

\subsection{Effective evaluation without a simulated execution}
